% ---------------------------------------------------------------------
\subsection{Equalizador BCD-MMSE}\label{subsec:bcd}
% ---------------------------------------------------------------------



O equalizador BCD-MMSE \textcolor{red}{(\textit{Block Coordinate Descent}--\textit{Minimum Mean Square Error}) \cite{zhao2025jsac}}, por sua vez, emprega atualizações distribuídas para tratar o acoplamento entre os blocos do sistema, mantendo a necessidade de definir quais informações devem ser compartilhadas entre as unidades.

No BCD-MMSE considerado neste trabalho, a estimativa local $\hat{\mathbf R}_{cc}$ participa da inicialização do equalizador e da construção da matriz utilizada nas atualizações da DU. Em particular, a matriz local a ser invertida é definida por:

\[
\mathbf G_c=
E_s\mathbf H_c\mathbf H_c^H+
\hat{\mathbf R}_{cc},
\]

em que $\mathbf H_c$ representa a parcela da matriz de canal associada ao cluster $c$. Dessa forma, a estimação e o condicionamento da covariância local afetam diretamente a viabilidade das operações matriciais necessárias à equalização.

% SEPARAÇÃO (2026-11-04): propriedades do algoritmo trazidas do rascunho (Partes 1.3 e 1.4)
%   e de METODOLOGIA/simulacao.tex.
\textcolor{red}{A DU $c$ usa $\hat{\mathbf R}_{cc}$ de duas formas. Na inicialização, usa a inversa $\hat{\mathbf R}_{cc}^{-1}$, pois o BCD-MMSE começa pela aproximação bloco-diagonal \cite{zhao2025jsac}. Nas iterações, usa $\mathbf G_c$, que soma $\hat{\mathbf R}_{cc}$ ao termo do canal \cite{zhao2025jsac}. Assim, $\mathbf G_c$ usa $\hat{\mathbf R}_{cc}$, e não a sua inversa. Além disso, $\mathbf G_c$ não depende da iteração \cite{zhao2025jsac}. Ela é formada e invertida uma vez, e a inversa é usada em todas as iterações.}
% FONTE: Algoritmo 3, paper.md:415 (W^0 pela Eq. 7); Eq. 7, paper.md:136-142 (R_cc^-1);
%   Eq. 30, paper.md:504 (G_c sem índice de iteração).

\textcolor{red}{O número $N$ de amostras também tem um custo de armazenamento e de comunicação. Cada DU guarda as $N$ amostras de ruído das suas antenas \cite{zhao2025jsac}. Trocar essas amostras entre as DUs exigiria transferir dados de tamanho $M \times N$ \cite{zhao2025jsac}. No BCD-MMSE, a banda de comunicação em cada iteração também depende de $N$ \cite{zhao2025jsac}. Como a estimativa só muda a cada bloco de coerência, esse custo é pago uma vez a cada $N_{\text{coh}}$ símbolos \cite{zhao2025jsac}.}
% FONTE: paper.md:493 (N_c em C^{M_c x N} na DU c); :130 (M x N); :471 (banda por
%   iteração depende de N); :561 e :571 (termos divididos por N_coh).


\subsubsection{Complexidade}\label{subsubsec:bcd-complexidade}
\textcolor{red}{No artigo base, o pré-processamento de cada DU tem complexidade $\mathcal{O}(M_c^3 + N M_c^2)$, e cada iteração do BCD-MMSE tem complexidade $\mathcal{O}(NMK + M M_c K)$ no total \cite{zhao2025jsac}. A Tabela~\ref{tab:bcd-complexidade} detalha o número de multiplicações complexas de cada parte do algoritmo em uma DU, por PRB e por $T_c$, para $K = 8$ e \sout{$N = 192$}.} \textcolor{red}{Como $N$ muda com a condição de teste, a tabela mostra o maior e o menor valor usados, $N = 192$ e $N = 24$ (Tabela~\ref{tab:condicoes_teste}). As linhas que não dependem de $N$ têm um só valor.}

\begin{table}[htbp]
\caption{\textcolor{red}{Multiplicações complexas por PRB e por $T_c$ em cada DU.}}
\label{tab:bcd-complexidade}
\centering
\footnotesize
\renewcommand{\arraystretch}{1.2}
\begin{tabular}{|p{3.6cm}|l|r|r|r|r|l|}
\hline
\raggedright Parte & Conta & \multicolumn{2}{c|}{$M_c = 16$} & \multicolumn{2}{c|}{$M_c = 8$} & Onde roda \tabularnewline
\cline{3-6}
 & & \textcolor{red}{$N = 192$} & \textcolor{red}{$N = 24$} & \textcolor{red}{$N = 192$} & \textcolor{red}{$N = 24$} & \tabularnewline
\hline
\raggedright Estimar $\hat{\mathbf R}_{cc}$ e aplicar Ledoit-Wolf & $N M_c^2 + N M_c + M_c^2$ & 52480 & \textcolor{red}{6784} & 13888 & \textcolor{red}{1792} & unidade de cálculo \tabularnewline
\raggedright Inicialização & $M_c^2 K + M_c K^2 + 2K^2 M_c + N K M_c$ & 29696 & \textcolor{red}{8192} & 14336 & \textcolor{red}{3584} & unidade de cálculo \tabularnewline
\raggedright Formar $\mathbf G_c$ & $M_c^2 K$ & \multicolumn{2}{c|}{2048} & \multicolumn{2}{c|}{512} & unidade de cálculo \tabularnewline
\raggedright Cada iteração & $2NKM_c + 3K^2 M_c + K M_c^2$ & 54272 & \textcolor{red}{11264} & 26624 & \textcolor{red}{5120} & unidade de cálculo \tabularnewline
\raggedright Inverter $\hat{\mathbf R}_{cc}$ & $M_c^3$ & \multicolumn{2}{c|}{4096} & \multicolumn{2}{c|}{512} & módulo 1 \tabularnewline
\raggedright Inverter $\mathbf G_c$ & $M_c^3$ & \multicolumn{2}{c|}{4096} & \multicolumn{2}{c|}{512} & módulo 2 \tabularnewline
\raggedright Inverter matriz $K \times K$ (só na DU $C$) & $K^3$ & \multicolumn{2}{c|}{512} & \multicolumn{2}{c|}{512} & módulo 3 \tabularnewline
\raggedright Filtro, por símbolo de dados & $12 K M_c$ & \multicolumn{2}{c|}{1536} & \multicolumn{2}{c|}{768} & unidade de cálculo \tabularnewline
\hline
\end{tabular}
\end{table}
% FONTE: software/Restrições - Parâmetros.csv, linhas 56-65 (modelo B). K = 8, N = 192: paper.md:577.
%        Colunas N = 24: mesmas fórmulas com N = 24 (menor N simulado, Tabela tab:n-condicao).
%        Estimar R_cc: paper.md:554 (O(M_c^3 + N M_c^2) por DU). Inicialização (X_c, H_c^H X_c, W_c^0,
%        A^0, b^0): paper.md:136-142 (Eq. 7) e :410-418. G_c: paper.md:504. Iteração (Eqs. 28-30):
%        paper.md:556 (O(NMK + M M_c K) no total), :443-449, :504. Filtro: paper.md:430-434.
% A EXPANDIR: explicar cada linha, o papel dos módulos 1-3 e o orçamento de tempo por PRB.
